% page 160 of the facsimile (image p-159 of the scan)
% block 160.0 x 237.7 mm ; left margin 8.0 mm
\pgc{-12.700mm}{0.000mm}{
\l{\cel{5}150.}
\l{}
\l{}
\l{\sou{exemplo}. I. Maniero esar, agar, di ulu, konsiderata kom im-}
\l{\cel{3}itinda. - II. La persono qua esas, agas, tale. - III.}
\l{\cel{3}To quon subisas ulu (desfortuno, puniseso, e c.), egard-}
\l{\cel{3}ata kom apta donar leciono. - IV. La persono qua subisas}
\l{\cel{3}ta desfortuno, ta puniso. - V. To quo eventis, konsider-}
\l{\cel{3}ata kom komparo-termino pri to simila quo povas eventar.}
\l{\cel{3}- VI. Fakto quan onu mencionas sekonde di aserto. -}
\l{\cel{3}VII. Extraktajo di texto quan onu citas, sekonde di expli-}
\l{\cel{3}ko gramatikala. - DEFIRS.}
\l{}
\l{\sou{exemplero}. I. Singla de la objekti formacita per tipo unika}
\l{\cel{3}quan onu reproduktis. - II. Singla de la specimeni di un}
\l{\cel{3}speco o di un varietato en kolektajo zoologiala, botanik-}
\l{\cel{3}ala, mineralogiala.\cel{3}DEFIRS.}
\l{}
\l{\sou{\textquotedbl{}exequatur\textquotedbl{}}. I. (\sou{Yuro-cienco}). Permiso legala necesa por ke}
\l{\cel{3}verdikto da arbitri, da tribunali, stranjera divenez exe-}
\l{\cel{3}kutebla. - II. Permiso quan bezonas aganto stranjera por}
\l{\cel{3}exekutar en la lando la misiono quan konfidis a lu lua}
\l{\cel{3}guvernistaro.}
\l{}
\l{\sou{exercar}. (\sou{trans}.)Formacar, fasonar (ulu) por igar lu prakti-}
\l{\cel{3}kar. - DEFIRS.}
\l{}
\l{\sou{exergo}. (\sou{numismatiko}). Spaco rezervita maxim-multa-kaze en}
\l{\cel{3}la infra parto di medalio, por recevar surskriburo (moto,}
\l{\cel{3}devizo, dato, e c.) - DEFIS.}
\l{}
\l{\sou{exfoliar}. (\sou{trans.}) Deprenar foliope, lamelope (ula parti di}
\l{\cel{3}substanco). - DEFIS.}
\l{}
\l{\sou{exhalar}. (\sou{trans.}) I. Emisar (gaso, vaporo, odoro). - II.}
\l{\cel{3}(\sou{pri animali, vejetanti}). Lasar eskapar tra la tisui ula}
\l{\cel{3}elementi eliminata de la organismo. - III. (\sou{metaf}.) Lasar}
\l{\cel{3}eskapar (sentimenti). - EFIS.}
\l{}
\l{\sou{exhaustar.}\cel{2}(\sou{trans.}) I. Sikigar per par-cherpo. - II. Vendar}
\l{\cel{3}la lasta exemplero (di libro, e c.). -III. Igar febla kom-}
\l{\cel{3}plete. - IV. (\sou{logiko}). Enumerar omna kazi, omna hipotezi po-}
\l{\cel{3}sibla di temo. - V. (\sou{matem}.)Uzar la procedo qua posibligas}
\l{\cel{3}neglijar la difero inter du grandori, per pruvar ke la di-}
\l{\cel{3}fero esas infinite mikra. - EFI.}
\l{}
\l{\sou{exhortar}. (\sou{trans., pri}) Probar igar (ulu) volar ulo, per dis-}
\l{\cel{3}kursi persuadiva. - EFIS.}
\l{}
\l{\sou{exilar}. (\sou{trans.}) Koaktar (ulu) a livar lua patrio, interdikt-}
\l{\cel{3}ante a lu retrovenor en olu. - DEFI.}
\l{}
\l{\sou{exino}. (\sou{biol.}) Strato extera, transformita a kutino (C6 H10 O)}
\l{\cel{3}e neegale dikeskinta (pinti e kresti, pori e plisuri),}
\l{\cel{3}di poleno-grano.}
}
